\section{Conclusion}
\label{sec:conclusion}

We have introduced \bfsg{}, the simplest redistricting algorithm in the
B-series platform.
Given $k$ seeds placed at maximally spread locations via BFS-farthest
selection, \bfsg{} grows $k$ districts simultaneously by repeatedly
assigning unassigned tracts to the adjacent district with the greatest
population deficit.
The algorithm requires zero hyperparameters beyond the standard compositor
inputs, runs in $O(n \log n)$, and guarantees contiguity by construction.

\paragraph{Position in the B-series algorithm landscape.}
\bfsg{} occupies the lowest-quality, fastest position in the structure-layer
algorithm spectrum.
The quality-runtime ordering across structure-layer algorithms is:
\begin{align*}
  \mathrm{PP:} &\quad \cvd{} \;\geq\; \bfsg{} \;>\; \metis{} \\
  \mathrm{EC:} &\quad \metis{} \;<\; \cvd{} \;\leq\; \bfsg{} \\
  \mathrm{Runtime:} &\quad \bfsg{} \;\approx\; \metis{} \;\ll\; \cvd{} \;\ll\; \sa{} \;\ll\; \be{}
\end{align*}
This ordering defines the operational role of each algorithm:
\begin{itemize}
\item \textbf{Fastest, zero-hyperparameter baseline}: \bfsg{}
  (recommended for geometric floor measurement and warm-start evaluation).
\item \textbf{EC-optimised single plan}: \metis{}
  (recommended for standard single-plan generation).
\item \textbf{PP-optimised, moderate cost}: \cvd{}
  (recommended when geographic compactness is the primary objective).
\item \textbf{High-quality optimised plan}: SA(steps=10) from B.19
  (recommended for publication-quality single-plan generation).
\item \textbf{Maximum quality, high cost}: \be{}$(T=100)$ or \cs{}
  (recommended for statutory map production).
\end{itemize}

\paragraph{Contiguity as a design feature.}
The contiguity guarantee (Proposition~\ref{prop:contiguity}) distinguishes
\bfsg{} from \metis{} and \cvd{}, both of which may produce disconnected
districts that require post-hoc contiguity repair.
\bfsg{}'s single-pass BFS growth is provably contiguous without any
repair step, making it the only structure-layer algorithm in the platform
with a formal contiguity guarantee.
This property is a legal asset: a BFS Growth plan can be submitted with
a proof that district contiguity was guaranteed algorithmically, not
merely verified post-hoc.

\paragraph{The geometric floor in context.}
The ``geometric floor'' framing is useful for legal and technical
communication.
If a submitted plan achieves higher PP than \bfsg{}, it demonstrates that
deliberate geographic attention was paid to compactness.
If a submitted plan achieves \emph{lower} PP than \bfsg{}, a zero-optimisation
greedy algorithm would have done better---a strong argument that the plan's
non-compactness is not an unavoidable consequence of geographic constraints.

\paragraph{Implementation.}
\bfsg{} is implemented in the \texttt{redist} binary under
\texttt{--structure bfs-growth}.
No additional parameters are required.
The YAML configuration is:
\begin{verbatim}
algorithm:
  structure: bfs-growth
  weights: geographic
  search: single
  balance_tolerance: 0.5
\end{verbatim}

\paragraph{Open questions.}
Several extensions are identified for future work:
\begin{enumerate}
\item \textbf{Composite priority function.} The current priority is
  current district population (deficit-only).
  A composite $w_{\mathrm{deficit}} \times \mathrm{deficit} + w_{\mathrm{dist}} \times d_{\mathrm{BFS}}(\mathrm{tract}, \mathrm{seed})$
  may improve PP at modest runtime cost.
  The spec notes this as a Phase~2 evaluation item.

\item \textbf{BFS Growth + ApportionRegions tree.}
  \ar{}~\citep{b11paper} links the bisection tree structure to the
  prime factorisation of $k$.
  Substituting \bfsg{} for \metis{} at each \ar{} node creates a
  zero-hyperparameter, prime-factor tree plan.
  This is a natural composition and requires no structural changes to
  the platform.

\item \textbf{Empirical warm-start validation.}
  Section~\ref{sec:warmstart} states the theoretical framework and
  expected directions for MCMC warm-start experiments.
  Executing the paired-chain autocorrelation experiments on NC, WI, and TX
  is the key empirical validation deferred from this paper.

\item \textbf{High-density states.}
  The spec notes that for high-density states (NY, CA, IL), seed quality
  matters more than the growth priority function.
  Evaluating whether BFS Growth's population-weighted seed~0 selection
  adequately handles dense urban graphs is a natural extension.
\end{enumerate}

These extensions inherit the three-layer compositor design and require no
structural changes to the codebase.
